\documentclass[11pt]{article}
\usepackage[utf8]{inputenc}
\usepackage[T1]{fontenc}
\usepackage{lmodern}
\usepackage[margin=1in]{geometry}
\usepackage{amsmath,amssymb,amsthm}
\usepackage{booktabs}
\usepackage{algorithm}
\usepackage{algorithmic}
\usepackage{microtype}
\usepackage{xcolor}
\usepackage{graphicx}
\usepackage[colorlinks=true,linkcolor=blue!50!black,citecolor=blue!50!black,urlcolor=blue!50!black]{hyperref}

\newtheorem{proposition}{Proposition}
\newtheorem{corollary}{Corollary}
\theoremstyle{remark}
\newtheorem{remark}{Remark}

\newcommand{\R}{\mathbb{R}}
\newcommand{\E}{\mathbb{E}}
\newcommand{\norm}[1]{\lVert #1\rVert}
\newcommand{\OPT}{\mathrm{OPT}}
\newcommand{\ceTwoAlpha}{0.5743}
\newcommand{\tauStar}{0.1927}
\newcommand{\ceTwoTrials}{20{,}000}
\newcommand{\ceTwoFailRecipe}{48.66\%}
\newcommand{\ceTwoFailRecipeLo}{47.96\%}
\newcommand{\ceTwoFailRecipeHi}{49.35\%}
\newcommand{\ceTwoZeroUpper}{0.019\%}
\newcommand{\ceTwoKppLloyd}{161.32}
\newcommand{\ceTwoKppShuffles}{2{,}000}
\newcommand{\tauLastFail}{0.1925}
\newcommand{\tauFirstOk}{0.193}
\newcommand{\tauStarHalf}{0.0964}
\newcommand{\failAtFive}{44.4\%}
\newcommand{\plateauFail}{14.37\%}
\newcommand{\plateauLo}{13.90\%}
\newcommand{\plateauHi}{14.87\%}
\newcommand{\failAtTen}{16.5\%}
\newcommand{\portTrials}{2{,}000}
\newcommand{\portRef}{49.00\%}
\newcommand{\portNumpy}{49.15\%}
\newcommand{\seedsRT}{10}
\newcommand{\BRT}{256}
\newcommand{\alphaRT}{0.402}
\newcommand{\ratioRTRecipe}{18.457}
\newcommand{\passRTRecipe}{2.0}
\newcommand{\winsRTRecipe}{0}
\newcommand{\fbRTRecipe}{0}
\newcommand{\timeRTRecipe}{0}
\newcommand{\convRTRecipe}{0}
\newcommand{\ratioRTPlus}{1.000}
\newcommand{\passRTPlus}{11.0}
\newcommand{\winsRTPlus}{0}
\newcommand{\fbRTPlus}{10}
\newcommand{\timeRTPlus}{0}
\newcommand{\convRTPlus}{10}
\newcommand{\ratioRTPlusS}{1.000}
\newcommand{\passRTPlusS}{20.0}
\newcommand{\winsRTPlusS}{0}
\newcommand{\fbRTPlusS}{10}
\newcommand{\timeRTPlusS}{0}
\newcommand{\convRTPlusS}{10}
\newcommand{\ratioRTRecipeC}{1.000}
\newcommand{\passRTRecipeC}{11.0}
\newcommand{\winsRTRecipeC}{0}
\newcommand{\fbRTRecipeC}{10}
\newcommand{\timeRTRecipeC}{0}
\newcommand{\convRTRecipeC}{10}
\newcommand{\ratioRTLloyd}{1.000}
\newcommand{\passRTLloyd}{7.0}
\newcommand{\winsRTLloyd}{0}
\newcommand{\fbRTLloyd}{0}
\newcommand{\timeRTLloyd}{0}
\newcommand{\convRTLloyd}{10}
\newcommand{\ratioRTLloydThree}{0.708}
\newcommand{\passRTLloydThree}{28.4}
\newcommand{\winsRTLloydThree}{7}
\newcommand{\fbRTLloydThree}{0}
\newcommand{\timeRTLloydThree}{0}
\newcommand{\convRTLloydThree}{0}
\newcommand{\ratioRTGuard}{0.954}
\newcommand{\passRTGuard}{14.9}
\newcommand{\winsRTGuard}{1}
\newcommand{\fbRTGuard}{7}
\newcommand{\timeRTGuard}{0}
\newcommand{\convRTGuard}{10}
\newcommand{\swapsRTGuard}{17}
\newcommand{\swapsDownRTGuard}{3}
\newcommand{\ratioRTGuardOnly}{1.000}
\newcommand{\passRTGuardOnly}{15.5}
\newcommand{\winsRTGuardOnly}{0}
\newcommand{\fbRTGuardOnly}{8}
\newcommand{\timeRTGuardOnly}{0}
\newcommand{\convRTGuardOnly}{10}
\newcommand{\swapsRTGuardOnly}{13}
\newcommand{\swapsDownRTGuardOnly}{1}
\newcommand{\ratioRTRecipeS}{1.000}
\newcommand{\passRTRecipeS}{17.9}
\newcommand{\winsRTRecipeS}{0}
\newcommand{\fbRTRecipeS}{10}
\newcommand{\timeRTRecipeS}{0}
\newcommand{\convRTRecipeS}{10}
\newcommand{\ratioRTNoRev}{0.997}
\newcommand{\passRTNoRev}{9.9}
\newcommand{\winsRTNoRev}{1}
\newcommand{\fbRTNoRev}{0}
\newcommand{\timeRTNoRev}{0}
\newcommand{\convRTNoRev}{10}
\newcommand{\ratioRTLloydConv}{1.000}
\newcommand{\passRTLloydConv}{7.0}
\newcommand{\winsRTLloydConv}{0}
\newcommand{\fbRTLloydConv}{0}
\newcommand{\timeRTLloydConv}{0}
\newcommand{\convRTLloydConv}{10}
\newcommand{\seedsRTB}{10}
\newcommand{\BRTB}{32}
\newcommand{\alphaRTB}{0.062}
\newcommand{\ratioRTBRecipe}{19.341}
\newcommand{\passRTBRecipe}{2.0}
\newcommand{\winsRTBRecipe}{0}
\newcommand{\fbRTBRecipe}{0}
\newcommand{\timeRTBRecipe}{0}
\newcommand{\convRTBRecipe}{0}
\newcommand{\ratioRTBPlus}{1.000}
\newcommand{\passRTBPlus}{11.0}
\newcommand{\winsRTBPlus}{0}
\newcommand{\fbRTBPlus}{10}
\newcommand{\timeRTBPlus}{0}
\newcommand{\convRTBPlus}{10}
\newcommand{\ratioRTBPlusS}{1.000}
\newcommand{\passRTBPlusS}{19.0}
\newcommand{\winsRTBPlusS}{0}
\newcommand{\fbRTBPlusS}{10}
\newcommand{\timeRTBPlusS}{0}
\newcommand{\convRTBPlusS}{10}
\newcommand{\ratioRTBRecipeC}{1.000}
\newcommand{\passRTBRecipeC}{11.0}
\newcommand{\winsRTBRecipeC}{0}
\newcommand{\fbRTBRecipeC}{10}
\newcommand{\timeRTBRecipeC}{0}
\newcommand{\convRTBRecipeC}{10}
\newcommand{\ratioRTBLloyd}{1.000}
\newcommand{\passRTBLloyd}{7.0}
\newcommand{\winsRTBLloyd}{0}
\newcommand{\fbRTBLloyd}{0}
\newcommand{\timeRTBLloyd}{0}
\newcommand{\convRTBLloyd}{10}
\newcommand{\ratioRTBLloydThree}{0.708}
\newcommand{\passRTBLloydThree}{28.4}
\newcommand{\winsRTBLloydThree}{7}
\newcommand{\fbRTBLloydThree}{0}
\newcommand{\timeRTBLloydThree}{0}
\newcommand{\convRTBLloydThree}{0}
\newcommand{\ratioRTBGuard}{0.981}
\newcommand{\passRTBGuard}{14.6}
\newcommand{\winsRTBGuard}{1}
\newcommand{\fbRTBGuard}{6}
\newcommand{\timeRTBGuard}{0}
\newcommand{\convRTBGuard}{10}
\newcommand{\swapsRTBGuard}{15}
\newcommand{\swapsDownRTBGuard}{5}
\newcommand{\ratioRTBGuardOnly}{1.000}
\newcommand{\passRTBGuardOnly}{17.8}
\newcommand{\winsRTBGuardOnly}{0}
\newcommand{\fbRTBGuardOnly}{10}
\newcommand{\timeRTBGuardOnly}{0}
\newcommand{\convRTBGuardOnly}{10}
\newcommand{\swapsRTBGuardOnly}{75}
\newcommand{\swapsDownRTBGuardOnly}{21}
\newcommand{\ratioRTBRecipeS}{1.000}
\newcommand{\passRTBRecipeS}{20.0}
\newcommand{\winsRTBRecipeS}{0}
\newcommand{\fbRTBRecipeS}{10}
\newcommand{\timeRTBRecipeS}{0}
\newcommand{\convRTBRecipeS}{10}
\newcommand{\ratioRTBNoRev}{0.997}
\newcommand{\passRTBNoRev}{10.9}
\newcommand{\winsRTBNoRev}{1}
\newcommand{\fbRTBNoRev}{0}
\newcommand{\timeRTBNoRev}{0}
\newcommand{\convRTBNoRev}{10}
\newcommand{\ratioRTBLloydConv}{1.000}
\newcommand{\passRTBLloydConv}{7.0}
\newcommand{\winsRTBLloydConv}{0}
\newcommand{\fbRTBLloydConv}{0}
\newcommand{\timeRTBLloydConv}{0}
\newcommand{\convRTBLloydConv}{10}
\newcommand{\seedsDTen}{10}
\newcommand{\BDTen}{64}
\newcommand{\alphaDTen}{0.384}
\newcommand{\ratioDTenRecipe}{1.016}
\newcommand{\passDTenRecipe}{2.0}
\newcommand{\winsDTenRecipe}{0}
\newcommand{\fbDTenRecipe}{0}
\newcommand{\timeDTenRecipe}{0}
\newcommand{\convDTenRecipe}{0}
\newcommand{\ratioDTenPlus}{1.008}
\newcommand{\passDTenPlus}{4.0}
\newcommand{\winsDTenPlus}{2}
\newcommand{\fbDTenPlus}{0}
\newcommand{\timeDTenPlus}{0}
\newcommand{\convDTenPlus}{0}
\newcommand{\ratioDTenPlusS}{1.005}
\newcommand{\passDTenPlusS}{10.1}
\newcommand{\winsDTenPlusS}{2}
\newcommand{\fbDTenPlusS}{0}
\newcommand{\timeDTenPlusS}{0}
\newcommand{\convDTenPlusS}{10}
\newcommand{\ratioDTenRecipeC}{1.016}
\newcommand{\passDTenRecipeC}{4.0}
\newcommand{\winsDTenRecipeC}{0}
\newcommand{\fbDTenRecipeC}{0}
\newcommand{\timeDTenRecipeC}{0}
\newcommand{\convDTenRecipeC}{0}
\newcommand{\ratioDTenLloyd}{1.000}
\newcommand{\passDTenLloyd}{17.6}
\newcommand{\winsDTenLloyd}{0}
\newcommand{\fbDTenLloyd}{0}
\newcommand{\timeDTenLloyd}{0}
\newcommand{\convDTenLloyd}{9}
\newcommand{\ratioDTenLloydThree}{0.994}
\newcommand{\passDTenLloydThree}{57.2}
\newcommand{\winsDTenLloydThree}{5}
\newcommand{\fbDTenLloydThree}{0}
\newcommand{\timeDTenLloydThree}{0}
\newcommand{\convDTenLloydThree}{0}
\newcommand{\ratioDTenGuard}{1.005}
\newcommand{\passDTenGuard}{10.3}
\newcommand{\winsDTenGuard}{2}
\newcommand{\fbDTenGuard}{0}
\newcommand{\timeDTenGuard}{0}
\newcommand{\convDTenGuard}{10}
\newcommand{\swapsDTenGuard}{1}
\newcommand{\swapsDownDTenGuard}{1}
\newcommand{\ratioDTenGuardOnly}{1.014}
\newcommand{\passDTenGuardOnly}{13.5}
\newcommand{\winsDTenGuardOnly}{0}
\newcommand{\fbDTenGuardOnly}{0}
\newcommand{\timeDTenGuardOnly}{0}
\newcommand{\convDTenGuardOnly}{10}
\newcommand{\swapsDTenGuardOnly}{1}
\newcommand{\swapsDownDTenGuardOnly}{1}
\newcommand{\ratioDTenRecipeS}{1.013}
\newcommand{\passDTenRecipeS}{12.5}
\newcommand{\winsDTenRecipeS}{0}
\newcommand{\fbDTenRecipeS}{0}
\newcommand{\timeDTenRecipeS}{0}
\newcommand{\convDTenRecipeS}{10}
\newcommand{\ratioDTenNoRev}{1.006}
\newcommand{\passDTenNoRev}{10.2}
\newcommand{\winsDTenNoRev}{2}
\newcommand{\fbDTenNoRev}{0}
\newcommand{\timeDTenNoRev}{0}
\newcommand{\convDTenNoRev}{10}
\newcommand{\ratioDTenLloydConv}{1.000}
\newcommand{\passDTenLloydConv}{17.9}
\newcommand{\winsDTenLloydConv}{0}
\newcommand{\fbDTenLloydConv}{0}
\newcommand{\timeDTenLloydConv}{0}
\newcommand{\convDTenLloydConv}{10}
\newcommand{\seedsDSixtyFour}{10}
\newcommand{\BDSixtyFour}{64}
\newcommand{\alphaDSixtyFour}{0.163}
\newcommand{\ratioDSixtyFourRecipe}{0.989}
\newcommand{\passDSixtyFourRecipe}{2.0}
\newcommand{\winsDSixtyFourRecipe}{9}
\newcommand{\fbDSixtyFourRecipe}{0}
\newcommand{\timeDSixtyFourRecipe}{0}
\newcommand{\convDSixtyFourRecipe}{0}
\newcommand{\ratioDSixtyFourPlus}{1.003}
\newcommand{\passDSixtyFourPlus}{4.0}
\newcommand{\winsDSixtyFourPlus}{2}
\newcommand{\fbDSixtyFourPlus}{0}
\newcommand{\timeDSixtyFourPlus}{0}
\newcommand{\convDSixtyFourPlus}{0}
\newcommand{\ratioDSixtyFourPlusS}{0.988}
\newcommand{\passDSixtyFourPlusS}{16.0}
\newcommand{\winsDSixtyFourPlusS}{9}
\newcommand{\fbDSixtyFourPlusS}{0}
\newcommand{\timeDSixtyFourPlusS}{0}
\newcommand{\convDSixtyFourPlusS}{10}
\newcommand{\ratioDSixtyFourRecipeC}{0.989}
\newcommand{\passDSixtyFourRecipeC}{4.0}
\newcommand{\winsDSixtyFourRecipeC}{9}
\newcommand{\fbDSixtyFourRecipeC}{0}
\newcommand{\timeDSixtyFourRecipeC}{0}
\newcommand{\convDSixtyFourRecipeC}{0}
\newcommand{\ratioDSixtyFourLloyd}{1.000}
\newcommand{\passDSixtyFourLloyd}{16.1}
\newcommand{\winsDSixtyFourLloyd}{0}
\newcommand{\fbDSixtyFourLloyd}{0}
\newcommand{\timeDSixtyFourLloyd}{0}
\newcommand{\convDSixtyFourLloyd}{10}
\newcommand{\ratioDSixtyFourLloydThree}{0.991}
\newcommand{\passDSixtyFourLloydThree}{52.4}
\newcommand{\winsDSixtyFourLloydThree}{7}
\newcommand{\fbDSixtyFourLloydThree}{0}
\newcommand{\timeDSixtyFourLloydThree}{0}
\newcommand{\convDSixtyFourLloydThree}{0}
\newcommand{\ratioDSixtyFourGuard}{1.000}
\newcommand{\passDSixtyFourGuard}{16.0}
\newcommand{\winsDSixtyFourGuard}{4}
\newcommand{\fbDSixtyFourGuard}{0}
\newcommand{\timeDSixtyFourGuard}{0}
\newcommand{\convDSixtyFourGuard}{10}
\newcommand{\swapsDSixtyFourGuard}{20}
\newcommand{\swapsDownDSixtyFourGuard}{15}
\newcommand{\ratioDSixtyFourGuardOnly}{0.994}
\newcommand{\passDSixtyFourGuardOnly}{13.1}
\newcommand{\winsDSixtyFourGuardOnly}{8}
\newcommand{\fbDSixtyFourGuardOnly}{0}
\newcommand{\timeDSixtyFourGuardOnly}{0}
\newcommand{\convDSixtyFourGuardOnly}{10}
\newcommand{\swapsDSixtyFourGuardOnly}{71}
\newcommand{\swapsDownDSixtyFourGuardOnly}{27}
\newcommand{\ratioDSixtyFourRecipeS}{0.982}
\newcommand{\passDSixtyFourRecipeS}{13.1}
\newcommand{\winsDSixtyFourRecipeS}{10}
\newcommand{\fbDSixtyFourRecipeS}{0}
\newcommand{\timeDSixtyFourRecipeS}{0}
\newcommand{\convDSixtyFourRecipeS}{10}
\newcommand{\ratioDSixtyFourNoRev}{1.004}
\newcommand{\passDSixtyFourNoRev}{13.5}
\newcommand{\winsDSixtyFourNoRev}{5}
\newcommand{\fbDSixtyFourNoRev}{0}
\newcommand{\timeDSixtyFourNoRev}{0}
\newcommand{\convDSixtyFourNoRev}{10}
\newcommand{\ratioDSixtyFourLloydConv}{1.000}
\newcommand{\passDSixtyFourLloydConv}{16.1}
\newcommand{\winsDSixtyFourLloydConv}{0}
\newcommand{\fbDSixtyFourLloydConv}{0}
\newcommand{\timeDSixtyFourLloydConv}{0}
\newcommand{\convDSixtyFourLloydConv}{10}
\newcommand{\seedsANN}{5}
\newcommand{\BANN}{2048}
\newcommand{\alphaANN}{0.437}
\newcommand{\ratioANNRecipe}{2.055}
\newcommand{\passANNRecipe}{2.0}
\newcommand{\winsANNRecipe}{0}
\newcommand{\fbANNRecipe}{0}
\newcommand{\timeANNRecipe}{0}
\newcommand{\convANNRecipe}{0}
\newcommand{\recOneANNRecipe}{0.735}
\newcommand{\recFourANNRecipe}{0.851}
\newcommand{\recSixteenANNRecipe}{0.961}
\newcommand{\ratioANNPlus}{1.000}
\newcommand{\passANNPlus}{97.4}
\newcommand{\winsANNPlus}{0}
\newcommand{\fbANNPlus}{5}
\newcommand{\timeANNPlus}{22}
\newcommand{\convANNPlus}{5}
\newcommand{\recOneANNPlus}{0.776}
\newcommand{\recFourANNPlus}{0.916}
\newcommand{\recSixteenANNPlus}{1.000}
\newcommand{\ratioANNPlusS}{1.334}
\newcommand{\passANNPlusS}{102.2}
\newcommand{\winsANNPlusS}{0}
\newcommand{\fbANNPlusS}{3}
\newcommand{\timeANNPlusS}{23}
\newcommand{\convANNPlusS}{5}
\newcommand{\recOneANNPlusS}{0.757}
\newcommand{\recFourANNPlusS}{0.891}
\newcommand{\recSixteenANNPlusS}{0.983}
\newcommand{\ratioANNRecipeC}{1.171}
\newcommand{\passANNRecipeC}{78.0}
\newcommand{\winsANNRecipeC}{0}
\newcommand{\fbANNRecipeC}{4}
\newcommand{\timeANNRecipeC}{18}
\newcommand{\convANNRecipeC}{4}
\newcommand{\recOneANNRecipeC}{0.768}
\newcommand{\recFourANNRecipeC}{0.907}
\newcommand{\recSixteenANNRecipeC}{0.992}
\newcommand{\ratioANNLloyd}{1.000}
\newcommand{\passANNLloyd}{26.0}
\newcommand{\winsANNLloyd}{0}
\newcommand{\fbANNLloyd}{0}
\newcommand{\timeANNLloyd}{6}
\newcommand{\convANNLloyd}{0}
\newcommand{\recOneANNLloyd}{0.775}
\newcommand{\recFourANNLloyd}{0.915}
\newcommand{\recSixteenANNLloyd}{1.000}
\newcommand{\ratioANNLloydThree}{0.972}
\newcommand{\passANNLloydThree}{78.0}
\newcommand{\winsANNLloydThree}{2}
\newcommand{\fbANNLloydThree}{0}
\newcommand{\timeANNLloydThree}{28}
\newcommand{\convANNLloydThree}{0}
\newcommand{\recOneANNLloydThree}{0.779}
\newcommand{\recFourANNLloydThree}{0.923}
\newcommand{\recSixteenANNLloydThree}{1.000}
\newcommand{\ratioANNGuard}{0.835}
\newcommand{\passANNGuard}{78.8}
\newcommand{\winsANNGuard}{5}
\newcommand{\fbANNGuard}{0}
\newcommand{\timeANNGuard}{28}
\newcommand{\convANNGuard}{5}
\newcommand{\swapsANNGuard}{145}
\newcommand{\swapsDownANNGuard}{120}
\newcommand{\recOneANNGuard}{0.829}
\newcommand{\recFourANNGuard}{0.951}
\newcommand{\recSixteenANNGuard}{1.000}
\newcommand{\ratioANNGuardOnly}{0.866}
\newcommand{\passANNGuardOnly}{76.0}
\newcommand{\winsANNGuardOnly}{5}
\newcommand{\fbANNGuardOnly}{0}
\newcommand{\timeANNGuardOnly}{28}
\newcommand{\convANNGuardOnly}{5}
\newcommand{\swapsANNGuardOnly}{148}
\newcommand{\swapsDownANNGuardOnly}{118}
\newcommand{\recOneANNGuardOnly}{0.825}
\newcommand{\recFourANNGuardOnly}{0.947}
\newcommand{\recSixteenANNGuardOnly}{1.000}
\newcommand{\ratioANNRecipeS}{1.352}
\newcommand{\passANNRecipeS}{106.8}
\newcommand{\winsANNRecipeS}{0}
\newcommand{\fbANNRecipeS}{3}
\newcommand{\timeANNRecipeS}{24}
\newcommand{\convANNRecipeS}{5}
\newcommand{\recOneANNRecipeS}{0.760}
\newcommand{\recFourANNRecipeS}{0.895}
\newcommand{\recSixteenANNRecipeS}{0.984}
\newcommand{\ratioANNLloydConv}{1.000}
\newcommand{\passANNLloydConv}{93.4}
\newcommand{\winsANNLloydConv}{0}
\newcommand{\fbANNLloydConv}{0}
\newcommand{\timeANNLloydConv}{22}
\newcommand{\convANNLloydConv}{5}
\newcommand{\recOneANNLloydConv}{0.776}
\newcommand{\recFourANNLloydConv}{0.916}
\newcommand{\recSixteenANNLloydConv}{1.000}
\newcommand{\seedsSIFT}{3}
\newcommand{\BSIFT}{32768}
\newcommand{\alphaSIFT}{0.728}
\newcommand{\ratioSIFTRecipe}{0.9998}
\newcommand{\passSIFTRecipe}{2.0}
\newcommand{\winsSIFTRecipe}{0}
\newcommand{\fbSIFTRecipe}{0}
\newcommand{\timeSIFTRecipe}{6}
\newcommand{\convSIFTRecipe}{0}
\newcommand{\recOneSIFTRecipe}{0.367}
\newcommand{\recFourSIFTRecipe}{0.691}
\newcommand{\recSixteenSIFTRecipe}{0.927}
\newcommand{\ratioSIFTPlus}{0.9998}
\newcommand{\passSIFTPlus}{4.0}
\newcommand{\winsSIFTPlus}{0}
\newcommand{\fbSIFTPlus}{0}
\newcommand{\timeSIFTPlus}{10}
\newcommand{\convSIFTPlus}{0}
\newcommand{\recOneSIFTPlus}{0.372}
\newcommand{\recFourSIFTPlus}{0.696}
\newcommand{\recSixteenSIFTPlus}{0.930}
\newcommand{\ratioSIFTPlusS}{0.9959}
\newcommand{\passSIFTPlusS}{305.0}
\newcommand{\winsSIFTPlusS}{3}
\newcommand{\fbSIFTPlusS}{0}
\newcommand{\timeSIFTPlusS}{801}
\newcommand{\convSIFTPlusS}{0}
\newcommand{\recOneSIFTPlusS}{0.372}
\newcommand{\recFourSIFTPlusS}{0.697}
\newcommand{\recSixteenSIFTPlusS}{0.929}
\newcommand{\ratioSIFTRecipeC}{0.9998}
\newcommand{\passSIFTRecipeC}{4.0}
\newcommand{\winsSIFTRecipeC}{0}
\newcommand{\fbSIFTRecipeC}{0}
\newcommand{\timeSIFTRecipeC}{10}
\newcommand{\convSIFTRecipeC}{0}
\newcommand{\recOneSIFTRecipeC}{0.367}
\newcommand{\recFourSIFTRecipeC}{0.691}
\newcommand{\recSixteenSIFTRecipeC}{0.927}
\newcommand{\ratioSIFTLloyd}{1.0000}
\newcommand{\passSIFTLloyd}{26.0}
\newcommand{\winsSIFTLloyd}{0}
\newcommand{\fbSIFTLloyd}{0}
\newcommand{\timeSIFTLloyd}{69}
\newcommand{\convSIFTLloyd}{0}
\newcommand{\recOneSIFTLloyd}{0.361}
\newcommand{\recFourSIFTLloyd}{0.688}
\newcommand{\recSixteenSIFTLloyd}{0.928}
\newcommand{\ratioSIFTLloydThree}{0.9999}
\newcommand{\passSIFTLloydThree}{78.0}
\newcommand{\winsSIFTLloydThree}{0}
\newcommand{\fbSIFTLloydThree}{0}
\newcommand{\timeSIFTLloydThree}{381}
\newcommand{\convSIFTLloydThree}{0}
\newcommand{\recOneSIFTLloydThree}{0.364}
\newcommand{\recFourSIFTLloydThree}{0.692}
\newcommand{\recSixteenSIFTLloydThree}{0.930}
\newcommand{\ratioSIFTGuard}{0.9966}
\newcommand{\passSIFTGuard}{301.0}
\newcommand{\winsSIFTGuard}{3}
\newcommand{\fbSIFTGuard}{0}
\newcommand{\timeSIFTGuard}{791}
\newcommand{\convSIFTGuard}{1}
\newcommand{\recOneSIFTGuard}{0.370}
\newcommand{\recFourSIFTGuard}{0.700}
\newcommand{\recSixteenSIFTGuard}{0.931}
\newcommand{\ratioSIFTLloydConv}{0.9962}
\newcommand{\passSIFTLloydConv}{301.0}
\newcommand{\winsSIFTLloydConv}{3}
\newcommand{\fbSIFTLloydConv}{0}
\newcommand{\timeSIFTLloydConv}{796}
\newcommand{\convSIFTLloydConv}{0}
\newcommand{\recOneSIFTLloydConv}{0.367}
\newcommand{\recFourSIFTLloydConv}{0.696}
\newcommand{\recSixteenSIFTLloydConv}{0.932}
\newcommand{\swapsSmallTotal}{213}
\newcommand{\swapsSmallDown}{74}
\newcommand{\swapsSmallDownPct}{35\%}
\newcommand{\swapsANNTotal}{293}
\newcommand{\swapsANNDown}{238}
\newcommand{\swapsANNDownPct}{81\%}
\newcommand{\tauSensFallbacks}{0}
\newcommand{\tauCostQuarter}{0.62\%}
\newcommand{\tauCostOne}{1.35\%}
\newcommand{\tauCostFour}{2.81\%}
\newcommand{\tailSeedMax}{0.3}
\newcommand{\tailMassMin}{33.8}
\newcommand{\tailOffRatio}{0.998}
\newcommand{\reservoirRT}{1.6}

\makeatletter\newcommand{\rows}[1]{\@@input #1 }\makeatother

\title{\textbf{Stabilizing Exponential Moving Average $k$-Means:\\ Repairing Two Counterexamples}}
\author{Rohan Arun\\ Super Powers AI\\ \texttt{rohan@getsupers.com}}
\date{November 2026}

\begin{document}
\maketitle

\begin{abstract}
e-means (Bessa et al., 2026) replaces Lloyd iterations with exponential moving averages of
cluster counts and sums and revives low-mass clusters. Assylbekov (2026) gave two
counterexamples: mass-based revival can destroy a global optimum, making the ratio of expected
losses to Lloyd unbounded under $k$-means++ seeding, and near-zero initial counts let a single
mini-batch move a decision boundary past an unvisited point, which happens with probability
$24/49$ in his second example. Working from the author's reference implementation, we separate
three candidate fixes and show which one resolves which counterexample. A \emph{seed-loss check}
(return Lloyd from the seeds whenever the e-means output is worse than the seeds) resolves the
first counterexample and, on every input, keeps the $8(\ln k+2)$ expected approximation bound of
$k$-means++, so the expected-loss ratio to Lloyd is bounded; it costs two data passes when it does
not trigger. A \emph{pseudo-count floor} $\tau$ resolves the second: a one-step argument gives the
threshold $\tau^\star=\tauStar$, and in $\ceTwoTrials$ shuffles per value the failure rate falls from
$\ceTwoFailRecipe$ to zero between $\tau=\tauLastFail$ and $\tau=\tauFirstOk$. Neither fix alone
resolves both. We call their combination e-means+. Off the counterexamples the picture is mixed.
On scikit-learn's digits with $k{=}64$, $\tau=1$ raises e-means' loss by \tauCostOne{}. A
cost-guarded revival rule we also test cuts loss to $\ratioANNGuard\times$ Lloyd's on a synthetic
heavy-tailed set, at the pass budget of three Lloyd restarts, but costs loss on digits with $k{=}64$. On SIFT1M, e-means+ matches the loss of 25 Lloyd iterations
($\ratioSIFTPlus\times$) in $\passSIFTPlus$ data passes instead of $\passSIFTLloyd$, and the check never triggers.
We report these costs alongside the fixes.
\end{abstract}

\section{Introduction}

e-means~\cite{bessa2026} is a mini-batch $k$-means method for large-scale vector quantization:
each cluster keeps exponential moving averages of its count and coordinate sum, a cosine-annealed
step size sets the averaging rate, and clusters whose averaged mass falls below a threshold are
revived by splitting the heaviest cluster. Its authors report reaching the loss of 25 Lloyd
iterations in about two epochs on large embedding sets~\cite{bessa2026}.

Assylbekov~\cite{assylbekov2026} gave two small one-dimensional counterexamples to the published
recipe, with a reference implementation~\cite{zh3nisrepo}.
\textbf{Counterexample 1 (revival)}: $m$ points at 1, $m$ at 3 and one at $D=m+4$, with $k=2$.
The singleton cluster holds $1/(2m+1)<5\%$ of the mass, revival overwrites it, and e-means ends at
$(1,4)$ with loss $m(m+1)$ from every $k$-means++ seed pair, while Lloyd's expected loss stays below
$6m$; the ratio of expected losses is therefore unbounded.
\textbf{Counterexample 2 (forgetting)}: $X=(0,0,15,29,29,29,29,29)$ with seeds $(15,29)$. With
initial counts $10^{-10}$, an early batch with a zero but without the point 15 can pull the left centre below 1, the point 15 changes sides and never returns, and e-means ends at the Lloyd
fixed point $(0,80/3)$ with loss $490/3$ instead of $150$, with probability exactly $24/49$.

Two Faiss pull requests already change the revival rule: a cost-based rule that ranks clusters by
their removal cost~\cite{pr5915} and an opt-in separation test~\cite{pr5908}. We ask a narrower
question: which minimal change resolves which counterexample, what does each change guarantee,
and what does it cost on data that is not adversarial?

\paragraph{Contributions.}
\begin{enumerate}
\item A component-level analysis on a line-for-line port of the reference code
  (Tables~\ref{tab:ce1}--\ref{tab:ce2}): a seed-loss check resolves Counterexample~1, a
  pseudo-count floor resolves Counterexample~2, and neither resolves the other.
\item A one-step bound on the pseudo-count (Proposition~\ref{prop:step}) whose value for
  Counterexample~2, $\tau^\star=\tauStar$, sits between the largest tested $\tau$ that still fails
  ($\tauLastFail$) and the smallest that never fails ($\tauFirstOk$) in our sweep.
\item A seed-loss check (Proposition~\ref{prop:check}) that keeps the $k$-means++ guarantee on every
  input and so bounds the expected-loss ratio to Lloyd. The guarantee is inherited from the
  seeding, not earned by e-means itself.
\item Measurements on synthetic and real data that show the costs of each change. Two results are
  negative: $\tau=1$ costs loss on digits, and on the small data sets a sample-based cost-guarded
  revival accepts more swaps that raise the full-data loss than lower it
  (\swapsSmallDown{} of \swapsSmallTotal{} lowered it). The same rule works well on a heavy-tailed
  many-cluster set.
\end{enumerate}
We ran one standard real benchmark, SIFT1M, on CPU with NumPy. We make no claim about speed at the scales e-means targets: we did not run Faiss, GPUs, or the larger benchmarks of~\cite{bessa2026}.

\section{The e-means recipe}\label{sec:recipe}

Let $X=\{x_1,\dots,x_n\}\subset\R^d$, $C=\{c_1,\dots,c_k\}$ and
$\Phi(X,C)=\sum_{x\in X}\min_j\norm{x-c_j}^2$. e-means keeps a count $N_j$ and sum $S_j$ per
cluster, initialised to $N_j=\tau$ and $S_j=\tau c_j^{(0)}$ with $\tau=10^{-10}$, and uses centres
$c_j=S_j/N_j$. Over $T$ epochs of shuffled batches of size $B$, step $t$ assigns the batch to the
nearest centres, giving per-cluster counts $n_j$ and sums $s_j$, and updates
\begin{equation}\label{eq:ema}
N_j\leftarrow(1-a_t)N_j+a_t n_j,\qquad S_j\leftarrow(1-a_t)S_j+a_t s_j,\qquad
a_t=\alpha_0\bigl(0.55+0.45\cos(\pi t/M)\bigr),
\end{equation}
with $M$ the total number of steps. After every fifth step, each cluster with
$N_j<(0.1/k)\sum_\ell N_\ell$ is revived: the heaviest cluster's centre $c$ and half its count
are given to both clusters, at $c(1\mp 2^{-10})$. The output is the mean of each cluster's
assignments during the last epoch. $B$ and $\alpha_0$ follow the printed laws
$B=\min\bigl(n,2^{\mathrm{round}(\log_2(0.035\,k^{0.13}n^{0.88}T^{0.81}))}\bigr)$ and
$\alpha_0=\min\bigl(1,0.824\,k^{-0.46}B^{0.90}n^{-0.45}T^{-0.11}\bigr)$. For $k=2$ these are the laws
implemented in~\cite{zh3nisrepo}; we read the factor $2$ there as $k$.

All one-dimensional results use our port of the reference code~\cite{zh3nisrepo}, which
reproduces its assertions (losses 110, 420 and 2550 for $m=10,20,50$; the trajectory and final
loss $490/3$ of Counterexample~2). The multi-dimensional results use a NumPy implementation of the
same recipe, which agrees with the port on Counterexample~1 and, over $\portTrials$ shuffles of
Counterexample~2, fails in $\portNumpy$ of runs against $\portRef$ for the port. It is not the Faiss \texttt{contrib/e\_means} code.

\section{Three candidate fixes}\label{sec:fixes}

\paragraph{C: seed-loss check.} Let $C_0$ be the seeds and $C_e$ the e-means output. Return $C_e$ if
$\Phi(X,C_e)\le\Phi(X,C_0)$, and otherwise return Lloyd run from $C_0$. $\Phi(X,C_0)$ comes for free
with $k$-means++, so the check costs one data pass when it does not trigger; we count two.

\paragraph{$\tau$: pseudo-count floor.} Initialise $N_j=\tau$ with $\tau=1$ instead of $10^{-10}$.
The seed then carries weight $\tau\prod_r(1-a_r)$ in $c_j$, which decays as data arrive.

\paragraph{G: cost-guarded revival.} A cluster below the same mass threshold is only a candidate,
and a tested candidate is not tested again for one epoch. A replacement
$x'$ is drawn by $D^2$ sampling from a uniform reservoir $R$ of $\min(n,\max(256,16k))$ points filled during the first epoch, and
the swap $c_j\leftarrow x'$ is accepted only if it lowers $\Phi(R,\cdot)$ by more than $10^{-3}\Phi(R,C)$.
This is the accept-if-improves swap of local search for $k$-means~\cite{kanungo2004} and of
LBG-U~\cite{fritzke1997}; the removal cost it uses is the quantity behind the cost-based revival
of~\cite{pr5915}. We do not claim it as new. Our one-dimensional port evaluates the test on all of
$X$, so there it is exact.

\paragraph{S: polish.} Run Lloyd from the e-means output until the assignment stops changing, then
apply C. We define \textbf{e-means+} as $\tau$ + C with the recipe's mass revival unchanged, and
report S and G as options (Algorithm~\ref{alg:plus}).

\begin{algorithm}[t]
\caption{e-means+ (options in brackets)}\label{alg:plus}
\begin{algorithmic}[1]
\REQUIRE $X$, seeds $C_0$ from $k$-means++ with $\Phi_0=\Phi(X,C_0)$, epochs $T$, $\tau=1$
\STATE $N_j\leftarrow\tau$, $S_j\leftarrow\tau c_{0,j}$ for all $j$
\STATE Run the recipe of Section~\ref{sec:recipe} [with revival G instead of mass revival]
\STATE $C\leftarrow$ last-epoch means [then $C\leftarrow$ Lloyd$(X,C)$ to convergence]
\IF{$\Phi(X,C)>\Phi_0$}
  \STATE $C\leftarrow$ Lloyd$(X,C_0)$ to convergence
\ENDIF
\RETURN $C$
\end{algorithmic}
\end{algorithm}

\section{What each fix guarantees}

\begin{proposition}[One-step boundary stability]\label{prop:step}
Suppose $x$ is strictly nearest to $c_j$, with gap $g=\min_{\ell\ne j}(\norm{x-c_\ell}-\norm{x-c_j})>0$.
If one update moves $c_j$ by $\delta_j$ and every other centre by at most $\varepsilon$, and
$\delta_j+\varepsilon<g$, then $x$ is still strictly nearest to $c_j$. Without revival, the update
\eqref{eq:ema} moves $c_j$ by
\begin{equation}\label{eq:disp}
\delta_j=\frac{a\,n_j}{(1-a)N_j+a\,n_j}\,\norm{\bar x_j-c_j},\qquad \bar x_j=s_j/n_j,
\end{equation}
so at the first step ($N_j=\tau$), if $\norm{\bar x_j-c_j}\le\rho$ and $\varepsilon<g$, the point stays
with $c_j$ whenever
\begin{equation}\label{eq:tau}
\tau>n_j\,\frac{a}{1-a}\Bigl(\frac{\rho}{g-\varepsilon}-1\Bigr).
\end{equation}
\end{proposition}
\begin{proof}
By the triangle inequality, $\norm{x-c'_j}\le\norm{x-c_j}+\delta_j<\norm{x-c_\ell}-\varepsilon\le\norm{x-c'_\ell}$
for every $\ell\ne j$. For \eqref{eq:disp}, write $S_j=N_jc_j$; then
$c'_j-c_j=\bigl((1-a)S_j+as_j-c_j((1-a)N_j+an_j)\bigr)/N'_j=an_j(\bar x_j-c_j)/N'_j$.
Requiring $a n_j\rho/((1-a)\tau+an_j)<g-\varepsilon$ and solving for $\tau$ gives \eqref{eq:tau}.
\end{proof}

\paragraph{Application to Counterexample~2.} With seeds $(15,29)$, the point 15 sits at the left
centre, so $g=14$. The worst first batch is $\{0,0,29,29\}$: $n_L=2$, $\bar x_L=0$, $\rho=15$, and the
right centre does not move ($\varepsilon=0$). With $a=\alpha_0=\ceTwoAlpha$ the bound is
$\tau^\star=2a/(14(1-a))=\tauStar$. In one dimension the condition is also necessary: the point 15
crosses at this step if and only if $\tau<\tau^\star$. The proposition covers the first step only.
The reference trajectory crosses at the third batch, and we have no proof for later steps. In the sweep of Table~\ref{tab:tau}, however, all failures stop between $\tau=\tauLastFail$ and
$\tau=\tauFirstOk$, which bracket $\tau^\star$.

\begin{remark}\label{rem:notau}
The bound grows linearly in the batch count $n_j$ and in $a/(1-a)$, so no fixed constant $\tau$
should be expected to prevent forgetting on every input. $\tau$ is a mitigation. The worst-case
guarantee comes from C.
\end{remark}

\begin{proposition}[Seed-loss check]\label{prop:check}
With C, $\Phi(X,C_{\mathrm{out}})\le\Phi(X,C_0)$ on every input and in every run.
If $C_0$ comes from $k$-means++~\cite{arthur2007}, then, since
$\E[\Phi(X,\mathrm{Lloyd}(C_0))]\ge\OPT$,
\[
\E[\Phi(X,C_{\mathrm{out}})]\le 8(\ln k+2)\,\OPT
\qquad\text{and}\qquad
\frac{\E[\Phi(X,C_{\mathrm{out}})]}{\E[\Phi(X,\mathrm{Lloyd}(C_0))]}\le 8(\ln k+2).
\]
\end{proposition}
\begin{proof}
If $C_e$ is kept, $\Phi(X,C_e)\le\Phi(X,C_0)$ by the test. Otherwise the output is Lloyd from
$C_0$, and Lloyd never increases $\Phi$ (a cluster that empties keeps its centre). The bounds
follow by taking expectations.
\end{proof}

This answers the unbounded-ratio statement of~\cite{assylbekov2026} for every input. It is a
safety net: it says nothing about when e-means beats its seeds.

\begin{proposition}[Polish]\label{prop:polish}
With S, the output is a Lloyd fixed point: each non-empty cluster's centre is the mean of the
points nearest to it.
\end{proposition}
\begin{proof}
Lloyd stops only when an assignment repeats, and its last update set each centre to the mean of
that assignment's cells. Lloyd terminates because $\Phi$ strictly decreases whenever the
assignment changes and there are finitely many assignments. The fallback also runs Lloyd to
convergence. Our implementation caps runs at 300 iterations. Every capped run converged on the
small data sets. On SIFT1M nearly all did not, so there the polished outputs are 300 Lloyd steps,
not certified fixed points (Section~\ref{sec:bench}).
\end{proof}

\paragraph{G gives no full-data guarantee.} An accepted swap lowers the reservoir loss
$\Phi(R,\cdot)$. It lowers $\Phi(X,\cdot)$ only when $R=X$, as in our one-dimensional runs. We
measured the full-data effect of every accepted swap in Section~\ref{sec:bench}.

\section{Results on the counterexamples}

\begin{table}[t]
\centering
\caption{Counterexample~1 (revival). Loss from the seeds $(1,D)$ and exact expected loss under
$k$-means++ (all seed pairs and their exact probabilities, as in~\cite{zh3nisrepo}). C: seed-loss
check; S: polish then check; G: cost-guarded revival. Batches are the whole data set, so runs
are deterministic.}\label{tab:ce1}
\begin{tabular}{lrrrrrr}
\toprule
 & \multicolumn{2}{c}{$m=10$} & \multicolumn{2}{c}{$m=20$} & \multicolumn{2}{c}{$m=50$}\\
Variant & seeds & $\E_{k\text{-means++}}$ & seeds & $\E_{k\text{-means++}}$ & seeds & $\E_{k\text{-means++}}$\\
\midrule
\rows{generated/tab_ce1.tex}
\bottomrule
\end{tabular}
\end{table}

\paragraph{Counterexample 1.} Table~\ref{tab:ce1}: the recipe and $\tau$ alone end at $(1,4)$ from
every seed pair. G alone keeps the singleton, because removing it costs $(D-2)^2$ while the best
split gains far less, and it then matches Lloyd from every seed pair. C alone also matches Lloyd:
from $(1,D)$ the recipe's loss $m(m+1)$ exceeds the seed loss $4m$, so the check falls back. From
$(1,3)$, Lloyd itself ends at $(1,4)$. The expected losses of every variant with C or G equal
Lloyd's exactly, so the expected-loss ratio is 1.

\begin{table}[t]
\centering
\caption{Counterexample~2 (forgetting) from seeds $(15,29)$, $B=4$, $T=25$. Prescribed:
loss on the reference trajectory. Failure: loss above 150, over $\ceTwoTrials$ shuffles with
95\% Wilson interval. Last column: expected loss under $k$-means++ (exact seed probabilities,
$\ceTwoKppShuffles$ shuffles per seed pair). Lloyd: 150 from the seeds and $\ceTwoKppLloyd$ under
$k$-means++.}\label{tab:ce2}
\begin{tabular}{lrrcrr}
\toprule
Variant & prescribed & failure & 95\% interval & mean loss & $\E_{k\text{-means++}}$\\
\midrule
\rows{generated/tab_ce2.tex}
\bottomrule
\end{tabular}
\end{table}

\paragraph{Counterexample 2.} Table~\ref{tab:ce2}: the recipe fails in $\ceTwoFailRecipe$ of
shuffles (95\% interval $\ceTwoFailRecipeLo$--$\ceTwoFailRecipeHi$, containing $24/49\approx48.98\%$).
G changes nothing, because no revival happens. C and S change nothing either: $(0,80/3)$ is a Lloyd
fixed point whose loss $490/3$ is below the seed loss 450. Every variant with $\tau=1$ never fails
(95\% upper bound $\ceTwoZeroUpper$).

\begin{table}[t]
\centering
\caption{Recipe with pseudo-count $\tau$ on Counterexample~2 ($\ceTwoTrials$ shuffles per value,
same shuffles for every $\tau$). Proposition~\ref{prop:step} gives $\tau^\star=\tauStar$.}\label{tab:tau}
\begin{tabular}{lrrr}
\toprule
$\tau$ & failure & 95\% upper & mean loss\\
\midrule
\rows{generated/tab_tau.tex}
\bottomrule
\end{tabular}
\end{table}

\paragraph{The $\tau$ sweep.} Table~\ref{tab:tau}: failures persist up to $\tau=\tauLastFail$ and
vanish from $\tau=\tauFirstOk$ on, bracketing $\tau^\star=\tauStar$. There is a second, earlier drop, from
\failAtFive{} at $\tau=0.05$ to \failAtTen{} at $\tau=0.1$. Proposition~\ref{prop:step} explains it: a first
batch with a single zero ($n_L=1$) crosses if and only if $\tau<\tau^\star/2=\tauStarHalf$.
Between $\tau^\star/2$ and $\tau^\star$, the only first batch that crosses is the one holding both
zeros and not the point 15, which occurs with probability $\binom52/\binom84=1/7\approx14.29\%$. Just
below $\tau^\star$, the failure rate is \plateauFail{} (95\% interval \plateauLo--\plateauHi), consistent
with that route alone. At smaller $\tau$, later-batch crossings add to the rate. At
$\tau=10^{-10}$, first-batch crossings account for $3/7\approx42.86\%$ of shuffles against
\ceTwoFailRecipe{} observed. The exact count in~\cite{zh3nisrepo} gives $2100/4900=3/7$ for the
first batch and $300/4900$ for the third, which sum to $24/49$. All failures stop inside the
interval that brackets $\tau^\star$.

\section{Beyond the counterexamples}\label{sec:bench}

\paragraph{Setup.} \emph{Rare tail}: $n=5{,}000$, $d=8$, $k=50$; 4{,}950 points around 45 centres
drawn uniformly in $[-5,5]^8$ ($\sigma=0.2$) and five clusters of 10 points ($0.2\%$ of the data
each) at $20+40e_i$ ($\sigma=0.05$). The laws give $B=\BRT$, $\alpha_0=\alphaRT$. A second run fixes
$B=32$ ($\alpha_0=\alphaRTB$). \emph{Digits}: the 1{,}797 $8\times8$ images of scikit-learn's
digits set~\cite{sklearn} ($d=64$), with $k=10$ and $k=64$ ($B=\BDTen$). \emph{Heavy-tailed ANN}: $n=50{,}000$,
$d=128$, 256 Gaussian modes ($\sigma=0.35$) with Zipf weights $r^{-1.15}$, the rarest 10\% of
modes pushed outward by a factor 3.5, $k=256$, $B=\BANN$, and 1{,}000 queries (half from tail
modes). All e-means variants use $T=2$ and the recipe's laws. Every method in a seed starts from
the same $k$-means++ seeds; there are $\seedsRT$ seeds per set ($\seedsANN$ for ANN). We report
the per-seed loss ratio to Lloyd-25 (mean $\pm$ standard deviation over seeds) and the number of
full data passes, counting epochs, Lloyd iterations and loss evaluations. We do not report
wall-clock time: our code is unoptimised NumPy on two CPU cores, unlike the GPU code
of~\cite{bessa2026}. ``Best of 3'' runs Lloyd-25 from the shared seeds and from two more
$k$-means++ draws.

\begin{table*}[t]
\centering
\caption{Loss relative to Lloyd-25 from the same seeds (mean $\pm$ s.d.\ of the per-seed ratio; lower
is better) and full data passes. C: seed-loss check; S: polish to convergence, then C; G:
cost-guarded revival.}\label{tab:bench}
\resizebox{\textwidth}{!}{\begin{tabular}{lrrrrrrrrrrrr}
\toprule
Method & \multicolumn{2}{c}{Rare tail} & \multicolumn{2}{c}{Rare tail, $B{=}32$} & \multicolumn{2}{c}{Digits, $k{=}10$} & \multicolumn{2}{c}{Digits, $k{=}64$} & \multicolumn{2}{c}{Heavy-tailed ANN} & \multicolumn{2}{c}{SIFT1M} \\
 & loss ratio & passes & loss ratio & passes & loss ratio & passes & loss ratio & passes & loss ratio & passes & loss ratio & passes \\
\midrule
Lloyd-25 & 1.000 {\scriptsize $\pm$0.000} & 7.0 & 1.000 {\scriptsize $\pm$0.000} & 7.0 & 1.000 {\scriptsize $\pm$0.000} & 17.6 & 1.000 {\scriptsize $\pm$0.000} & 16.1 & 1.000 {\scriptsize $\pm$0.000} & 26.0 & 1.000 {\scriptsize $\pm$0.000} & 26.0 \\
Lloyd-25, best of 3 & 0.708 {\scriptsize $\pm$0.257} & 28.4 & 0.708 {\scriptsize $\pm$0.257} & 28.4 & 0.994 {\scriptsize $\pm$0.013} & 57.2 & 0.991 {\scriptsize $\pm$0.009} & 52.4 & 0.972 {\scriptsize $\pm$0.040} & 78.0 & 1.000 {\scriptsize $\pm$0.000} & 78.0 \\
e-means (recipe) & 18.457 {\scriptsize $\pm$10.760} & 2.0 & 19.341 {\scriptsize $\pm$10.967} & 2.0 & 1.016 {\scriptsize $\pm$0.014} & 2.0 & 0.989 {\scriptsize $\pm$0.008} & 2.0 & 2.055 {\scriptsize $\pm$0.168} & 2.0 & 1.000 {\scriptsize $\pm$0.000} & 2.0 \\
e-means + C & 1.000 {\scriptsize $\pm$0.000} & 11.0 & 1.000 {\scriptsize $\pm$0.000} & 11.0 & 1.016 {\scriptsize $\pm$0.014} & 4.0 & 0.989 {\scriptsize $\pm$0.008} & 4.0 & 1.171 {\scriptsize $\pm$0.382} & 78.0 & 1.000 {\scriptsize $\pm$0.000} & 4.0 \\
\textbf{e-means+} ($\tau$ + C) & 1.000 {\scriptsize $\pm$0.000} & 11.0 & 1.000 {\scriptsize $\pm$0.000} & 11.0 & 1.008 {\scriptsize $\pm$0.011} & 4.0 & 1.003 {\scriptsize $\pm$0.010} & 4.0 & 1.000 {\scriptsize $\pm$0.000} & 97.4 & 1.000 {\scriptsize $\pm$0.000} & 4.0 \\
e-means + S & 1.000 {\scriptsize $\pm$0.000} & 17.9 & 1.000 {\scriptsize $\pm$0.000} & 20.0 & 1.013 {\scriptsize $\pm$0.012} & 12.5 & 0.982 {\scriptsize $\pm$0.007} & 13.1 & 1.352 {\scriptsize $\pm$0.483} & 106.8 & -- & -- \\
$\tau$ + S & 1.000 {\scriptsize $\pm$0.000} & 20.0 & 1.000 {\scriptsize $\pm$0.000} & 19.0 & 1.005 {\scriptsize $\pm$0.011} & 10.1 & 0.988 {\scriptsize $\pm$0.009} & 16.0 & 1.334 {\scriptsize $\pm$0.458} & 102.2 & 0.996 {\scriptsize $\pm$0.000} & 305.0 \\
Guard + S & 1.000 {\scriptsize $\pm$0.000} & 15.5 & 1.000 {\scriptsize $\pm$0.000} & 17.8 & 1.014 {\scriptsize $\pm$0.010} & 13.5 & 0.994 {\scriptsize $\pm$0.007} & 13.1 & 0.866 {\scriptsize $\pm$0.062} & 76.0 & -- & -- \\
$\tau$ + guard + S & 0.954 {\scriptsize $\pm$0.145} & 14.9 & 0.981 {\scriptsize $\pm$0.061} & 14.6 & 1.005 {\scriptsize $\pm$0.011} & 10.3 & 1.000 {\scriptsize $\pm$0.010} & 16.0 & 0.835 {\scriptsize $\pm$0.042} & 78.8 & 0.997 {\scriptsize $\pm$0.000} & 301.0 \\
\bottomrule
\end{tabular}
}
\end{table*}

\begin{table}[t]
\centering
\small
\caption{Per-seed outcome against Lloyd-25: wins / losses (loss more than 0.1\% lower / higher) /
times the seed-loss check fell back to Lloyd.}\label{tab:wins}
\resizebox{\textwidth}{!}{\begin{tabular}{lcccccc}
\toprule
Variant & Rare tail & Rare tail, $B{=}32$ & Digits, $k{=}10$ & Digits, $k{=}64$ & Heavy-tailed ANN & SIFT1M \\
\midrule
e-means (recipe) & 0/10/0 & 0/10/0 & 0/8/0 & 9/0/0 & 0/5/0 & 0/0/0 \\
e-means + C & 0/0/10 & 0/0/10 & 0/8/0 & 9/0/0 & 0/1/4 & 0/0/0 \\
\textbf{e-means+} ($\tau$ + C) & 0/0/10 & 0/0/10 & 2/6/0 & 2/6/0 & 0/0/5 & 0/0/0 \\
e-means + S & 0/0/10 & 0/0/10 & 0/7/0 & 10/0/0 & 0/2/3 & -- \\
$\tau$ + S & 0/0/10 & 0/0/10 & 2/5/0 & 9/0/0 & 0/2/3 & 3/0/0 \\
Guard + S & 0/0/8 & 0/0/10 & 0/8/0 & 8/1/0 & 5/0/0 & -- \\
$\tau$ + guard + S & 1/0/7 & 1/0/6 & 2/5/0 & 4/4/0 & 5/0/0 & 3/0/0 \\
\bottomrule
\end{tabular}
}
\end{table}

\paragraph{The check prevents the rare-tail collapse.} Mass revival destroys tail clusters. The
$k$-means++ seeds cover all five (every tail point lies within \tailSeedMax{} of a seed), but after the
recipe, in every seed, some tail point is at least \tailMassMin{} from every centre. The recipe's
loss is on average $\ratioRTRecipe\times$ Lloyd's ($\ratioRTBRecipe\times$ with $B=32$); with revival
turned off it is $\tailOffRatio\times$. The check falls back in all $\seedsRT$ seeds (Table~\ref{tab:wins}), so e-means+ equals
Lloyd there, at $\passRTPlus$ passes instead of Lloyd's $\passRTLloyd$. This is the
Counterexample~1 mechanism at a larger scale, and the check handles it the same way. $\tau$ + G + S ends slightly below Lloyd on average ($\ratioRTGuard$, within one standard
deviation), and it falls back in $\fbRTGuard$ of the $\seedsRT$ seeds.

\paragraph{Digits.} With $k=10$, the recipe is worse than Lloyd-25 ($\ratioDTenRecipe$) and $\tau$
helps ($\ratioDTenPlus$), at $\passDTenPlus$ passes against Lloyd's $\passDTenLloyd$. With $k=64$,
the recipe ends below Lloyd-25 on every seed ($\ratioDSixtyFourRecipe$; by more than 0.1\% on
$\winsDSixtyFourRecipe$ of $\seedsRT$), but $\tau=1$ removes that advantage ($\ratioDSixtyFourPlus$). Table~\ref{tab:tausens}
shows that the cost grows with $\tau$. Here the batch holds about one point per cluster, so the
pseudo-count is not small next to the data. With $k=64$, G also costs loss: G + S ends at
$\ratioDSixtyFourGuardOnly$ against $\ratioDSixtyFourRecipeS$ for the recipe + S. With $k=10$, the
two are within noise. Turning revival off ($\tau$ + S without revival, in the results files)
ends at \ratioDSixtyFourNoRev{} against \ratioDSixtyFourPlusS{} with $k=64$, so the recipe's mass
revival is useful on this data.

\begin{table}[t]
\centering
\small
\caption{Change in final loss of the recipe + C on digits as $\tau$ varies, relative to
$\tau=10^{-10}$ (mean over the same 10 seeds). The check never fell back.}\label{tab:tausens}
\begin{tabular}{lrrrrr}
\toprule
$\tau$ & $10^{-10}$ & 0.05 & 0.25 & 1 & 4 \\
\midrule
Digits, $k{=}10$ & +0.00\% & +0.09\% & -0.02\% & -0.75\% & -0.31\% \\
Digits, $k{=}64$ & +0.00\% & +0.13\% & +0.62\% & +1.35\% & +2.81\% \\
\bottomrule
\end{tabular}

\end{table}

\begin{table}[t]
\centering
\small
\caption{Heavy-tailed ANN ($n=50{,}000$, $d=128$, $k=256$, $\seedsANN$ seeds): loss relative to
Lloyd-25, IVF-Flat Recall@10 at 1, 4 and 16 probes (mean, with s.d.\ for one probe), and data
passes.}\label{tab:ann}
\begin{tabular}{lrrrrr}
\toprule
Variant & loss ratio & R@10, 1 probe & 4 probes & 16 probes & passes\\
\midrule
\rows{generated/tab_ann.tex}
\bottomrule
\end{tabular}
\end{table}

\paragraph{Heavy-tailed ANN.} Table~\ref{tab:ann}. Here mass revival is costly: the recipe ends at
$\ratioANNRecipe\times$ Lloyd's loss, and polishing does not repair it, because the polished
result sits in a poor Lloyd optimum. The recipe + C falls back in $\fbANNRecipeC$ of the
$\seedsANN$ seeds, and $\tau$ + S in $\fbANNPlusS$, ending far above Lloyd in the others.
e-means+ falls back in all $\seedsANN$ seeds, so it returns Lloyd's result (run to convergence)
at $\passANNPlus$ passes against $\passANNLloyd$ for Lloyd-25. Its value here is the guarantee
alone. The cost-guarded variants do best: $\tau$ + G + S reaches $\ratioANNGuard\times$ Lloyd's
loss, beating Lloyd-25 on every seed. It also beats the best of three Lloyd runs
($\ratioANNLloydThree$) at a similar budget ($\passANNGuard$ against $\passANNLloydThree$ passes),
and raises one-probe Recall@10 from $\recOneANNLloyd$ to $\recOneANNGuard$. Of its accepted swaps,
$\swapsDownANNGuard$ of $\swapsANNGuard$ lowered the full-data loss. With 256 modes of very
unequal weight, a $D^2$-sampled swap that moves a centre from a crowded mode to an uncovered tail
mode is usually a real improvement. Lloyd cannot make that move.


\IfFileExists{generated/sift_section.tex}{\paragraph{SIFT1M.} We trained on 1{,}000,000 SIFT descriptors ($d=128$) with $k=1024$,
$T=2$, $B=32768$ and $\alpha_0=0.728$ from the recipe's laws, over 3 seeds,
and evaluated IVF-Flat Recall@10 on the full base with the data set's ground truth
(Table~\ref{tab:sift}). G was run without the per-swap audit.

\begin{table}[t]
\centering
\small
\caption{SIFT1M: loss relative to Lloyd-25 (mean $\pm$ s.d.\ of the per-seed ratio), Recall@10 at 1, 4
and 16 probes, data passes, and seeds in which the seed-loss check fell back.}\label{tab:sift}
\begin{tabular}{lrrrrrr}
\toprule
Variant & loss ratio & R@10, 1 & 4 & 16 & passes & fell back\\
\midrule
Lloyd-25 & 1.0000 {\scriptsize $\pm$0.0000} & 0.361 & 0.688 & 0.928 & 26.0 & 0 \\
Lloyd-25, best of 3 & 0.9999 {\scriptsize $\pm$0.0002} & 0.364 & 0.692 & 0.930 & 78.0 & 0 \\
e-means (recipe) & 0.9998 {\scriptsize $\pm$0.0001} & 0.367 & 0.691 & 0.927 & 2.0 & 0 \\
e-means + C & 0.9998 {\scriptsize $\pm$0.0001} & 0.367 & 0.691 & 0.927 & 4.0 & 0 \\
\textbf{e-means+} ($\tau$ + C) & 0.9998 {\scriptsize $\pm$0.0002} & 0.372 & 0.696 & 0.930 & 4.0 & 0 \\
$\tau$ + S & 0.9959 {\scriptsize $\pm$0.0003} & 0.372 & 0.697 & 0.929 & 305.0 & 0 \\
$\tau$ + guard + S & 0.9966 {\scriptsize $\pm$0.0003} & 0.370 & 0.700 & 0.931 & 301.0 & 0 \\
\bottomrule
\end{tabular}
\end{table}
}{}
\IfFileExists{generated/sift_section.tex}{\paragraph{SIFT1M results.} On this real benchmark the fixes cost nothing. The recipe reaches
$\ratioSIFTRecipe\times$ the loss of Lloyd-25 in $\passSIFTRecipe$ passes, and e-means+ reaches
$\ratioSIFTPlus\times$ in $\passSIFTPlus$. The check fell back in none of the $\seedsSIFT$ seeds, and
$\tau=1$ changes the loss by less than $0.01\%$. Because every method here runs in the same NumPy
code on the same machine (a MacBook Pro CPU), we also give wall-clock times: e-means+
took $\timeSIFTPlus$\,s against $\timeSIFTLloyd$\,s for Lloyd-25. These are not comparable to the
GPU timings of~\cite{bessa2026}. One-probe Recall@10 is $\recOneSIFTPlus$ for e-means+ and
$\recOneSIFTLloyd$ for Lloyd-25, and every variant lies within 1.5 points. Lloyd run for 300
iterations lowers the loss by about $0.4\%$ ($\ratioSIFTLloydConv\times$) at $\passSIFTLloydConv$
passes, and the polished variants match it ($\tau$ + S: $\ratioSIFTPlusS\times$). These runs
hit the 300-iteration cap without converging ($\convSIFTLloydConv$ of $\seedsSIFT$ Lloyd runs and
$\convSIFTPlusS$ of $\seedsSIFT$ $\tau$ + S runs converged), so S does not deliver a fixed point at
this scale within our cap. G adds nothing here ($\ratioSIFTGuard\times$ against $\ratioSIFTPlusS\times$
for $\tau$ + S). The guard's gain on the heavy-tailed synthetic set does not carry over to SIFT1M.
}{}

\paragraph{Swap audit.} On the rare-tail and digits runs with G, \swapsSmallDown{} of
\swapsSmallTotal{} accepted swaps (\swapsSmallDownPct) lowered the full-data loss at the moment they
were made. On the heavy-tailed set, \swapsANNDown{} of \swapsANNTotal{} did (\swapsANNDownPct).
On the rare-tail set the reservoir holds 800 of the 5{,}000 points, so a 10-point tail cluster appears
\reservoirRT{} times on average, and its estimated removal cost is noisy. The reservoir test is a heuristic. Where G helps,
the help comes from many good swaps, not from a guarantee.

\section{Related work}
Mini-batch $k$-means~\cite{sculley2010} averages with step $1/\text{count}$. e-means~\cite{bessa2026}
uses annealed exponential averages and revival. Fritzke's LBG-U~\cite{fritzke1997} moves the
centre with the lowest utility (removal cost) and keeps the move only if the error drops after
re-running $k$-means. Single-swap local search with the same acceptance rule is a constant-factor
approximation~\cite{kanungo2004}. Faiss PR~\#5915~\cite{pr5915} makes a utility-based revival the
default in e-means. PR~\#5908~\cite{pr5908} adds an opt-in separation test. The seed-loss check
relies only on Lloyd's monotonicity and the $k$-means++ bound~\cite{arthur2007}.

\section{Limitations}
(i) Our only large real benchmark is SIFT1M, run with NumPy on CPU and few seeds; we did not run Faiss, GPUs, or the larger data sets of~\cite{bessa2026} (e.g.\ DEEP). (ii) The multi-dimensional results come from a NumPy re-implementation of the recipe.
(iii) Proposition~\ref{prop:step} covers one step. The full sweep result in Table~\ref{tab:tau} is
empirical (95\% upper bound $\ceTwoZeroUpper$ per $\tau$). (iv) $\tau=1$ was fixed before the
real-data runs. Table~\ref{tab:tausens} suggests a smaller floor, or one scaled to the expected
per-cluster batch count, but we have not tested such a rule. (v) The data sets are small and the
seed counts modest, so differences within one standard deviation should not be read as wins.

\section{Conclusion}
The two counterexamples have separate causes and separate fixes. A two-pass seed-loss check
resolves the revival counterexample and bounds the expected-loss ratio to Lloyd on every input.
A pseudo-count floor above a computable one-step threshold resolves the forgetting
counterexample. Off the counterexamples, the floor can cost loss, and a sample-based
accept-if-improves revival depends on the data: it helps on a heavy-tailed many-cluster set and
hurts on digits. The check costs at most one Lloyd run from the seeds, the run e-means was meant
to replace, and nothing beyond two passes when it does not trigger. It is the one change we
recommend in every setting we tested.

\paragraph{Reproducibility.} Code, results and the scripts that generate every table and number in
this paper: \texttt{emeans-plus} (\texttt{benchmarks/run\_experiments.py},
\texttt{benchmarks/make\_tables.py}, \texttt{tests/}).

\begin{thebibliography}{10}
\bibitem{arthur2007} D.~Arthur and S.~Vassilvitskii. $k$-means++: the advantages of careful seeding.
  In \emph{Proc.\ SODA}, pages 1027--1035, 2007.
\bibitem{assylbekov2026} Z.~Assylbekov. Two counterexamples for e-means. arXiv:2611.22065, 2026.
\bibitem{bessa2026} S.~Bessa, P.-E.~Mazar\'e, G.~Szilvasy, H.~Touvron, D.~Picard, M.~Douze, and H.~J\'egou.
  e-means: outperforming $k$-means (in one-tenth the epochs). arXiv:2611.09501, 2026.
\bibitem{fritzke1997} B.~Fritzke. The LBG-U method for vector quantization -- an improvement over LBG
  inspired from neural networks. \emph{Neural Processing Letters}, 5(1):35--45, 1997.
\bibitem{kanungo2004} T.~Kanungo, D.~M.~Mount, N.~S.~Netanyahu, C.~D.~Piatko, R.~Silverman, and A.~Y.~Wu.
  A local search approximation algorithm for $k$-means clustering. \emph{Computational Geometry},
  28(2--3):89--112, 2004.
\bibitem{sklearn} F.~Pedregosa et al. Scikit-learn: machine learning in Python.
  \emph{Journal of Machine Learning Research}, 12:2825--2830, 2011.
\bibitem{pr5908} rzr1331. e-means: opt-in revival=``auto''/``separation'' that keeps real small
  clusters (NumpyEmeans). \url{https://github.com/facebookresearch/faiss/pull/5908}, 2026.
\bibitem{pr5915} lcabannes. e-means: cost-based cluster revival.
  \url{https://github.com/facebookresearch/faiss/pull/5915}, 2026.
\bibitem{sculley2010} D.~Sculley. Web-scale $k$-means clustering. In \emph{Proc.\ WWW}, pages 1177--1178, 2010.
\bibitem{zh3nisrepo} Z.~Assylbekov. Reference code for ``Two counterexamples for e-means''.
  \url{https://github.com/zh3nis/e-means-counterexamples}, 2026.
\end{thebibliography}
\end{document}
